\section{Related Work}
\label{sec:related}

\para{RLHF and preference optimization.}
The standard RLHF pipeline trains a reward model on pairwise human comparisons and optimizes a policy via PPO \citep{ouyang2022training, bai2022training, stiennon2020learning}. DPO \citep{rafailov2023direct} simplifies this by directly optimizing the policy with a binary cross-entropy loss, eliminating the need for a separate reward model while remaining mathematically equivalent under Bradley-Terry assumptions. Both approaches aggregate preferences across annotators, which \citet{chakraborty2024maxmin} show is provably insufficient to represent diverse preferences---a single reward model cannot capture heterogeneous annotator values. Our work uses DPO as a controlled testbed to isolate the effect of training on individual versus aggregated preferences.

\para{Length bias and verbosity in RLHF.}
Several studies have identified verbosity as a systematic failure mode of RLHF. \citet{shen2023loose} show that reward models learn a spurious length-to-reward correlation from human annotation data, which PPO then exploits. They propose a Product-of-Experts approach that reduces output length by approximately 15\% while maintaining quality. \citet{saito2023verbosity} quantify verbosity bias in both human annotators and LLM judges, finding that GPT-4 exhibits a verbosity bias score of 0.328 and that humans in standard datasets also tend to prefer longer responses. Unlike these works, which focus on mitigating length bias within multi-annotator RLHF, we ask whether the bias \emph{disappears} when training on a single annotator whose preferences are known.

\para{Sycophancy and reward hacking.}
\citet{sharma2024towards} provide a comprehensive analysis of sycophancy in language models, showing that it is present \emph{before} RL training and is amplified during optimization. Overoptimization of imperfect reward models is a related concern: \citet{gao2023scaling} characterize how reward model quality degrades with optimization pressure, and \citet{coste2023reward} show that ensembling reward models can mitigate this effect. These findings suggest that even a single-annotator reward model would be vulnerable to overoptimization, producing outputs that exploit the annotator's biases rather than genuinely reflecting their preferences.

\para{Personalized and pluralistic alignment.}
A growing body of work addresses the heterogeneity of human preferences. \citet{xie2025survey} survey personalized alignment methods, categorizing approaches as training-time (user-specific parameters, steerable models), inference-time (prompting, guided decoding), and user-modeling-based. \citet{park2024rlhf} provide a formal framework with sample complexity guarantees for personalized versus aggregated RLHF. \citet{li2024preference} show that KL-divergence regularization causes ``preference collapse'' where minority preferences are disregarded. Most personalized alignment work uses synthetic preferences; \prism \citep{kirk2024prism} is one of the few datasets with real per-user identifiers at scale. Our work differs from these methods-focused papers by using single-annotator training as a \emph{diagnostic tool} to attribute verbosity and blandness to specific pipeline stages, rather than proposing a new personalization method.
